\section{Symmetric profiles and mixed Gamma ratios}
\label{ag:sec:profiles}

The affine argument does not depend on the particular uniform interval
average in the Gamma ratio. Let $K(z)\in\C[[z^2]]$ be even, with
$K(0)=1$, and form
\begin{equation}\label{ag:eq:profile-density}
 \mathcal C_K(t)=\exp\bigl(tK(t\D)u_0\bigr).
\end{equation}
Its coefficients are differential polynomials, and its normal form again
has an Appell factor. Denote the coefficient of $u_r^2t^{2r+2}$ after
removing that factor by $g_{K,r}(y)$, with $y=u_1t^2/2$.

\begin{theorem}[The symmetric-profile coefficient theorem]
\label{ag:thm:profile}
For every $r\ge2$,
\begin{equation}\label{ag:eq:profile-symbol}
 \boxed{
 g_{K,r}(y)=[w^r]\frac12
 K\left(\sqrt{y^2-w}-y\right)
 K\left(\sqrt{y^2-w}+y\right).}
\end{equation}
The product is even in the square-root variable, so its interpretation
as a formal series in $y,w$ is unambiguous. Formula
\eqref{ag:eq:profile-symbol} determines all coefficients of
$u_0^a u_1^k u_r^2$ in $\NF\mathcal C_K$.
\end{theorem}
\begin{proof}
Under \eqref{ag:eq:affine}, the quadratic term is
$t^2W(K(t\D)v)^2/2$. Its symbol is
$t^2K(h-y)K(h+y)/2$, where $h^2=y^2-t^2\eta$.
The primitive has Appell covariance by
Lemma~\ref{ag:lem:appell-primitive}, and normal support excludes any
other contribution to $\eta^r$ for $r\ge2$. Comparing coefficients
of $t^{2r+2}\eta^r$ proves the formula. Evenness of $K$ makes the
product invariant under $h\mapsto-h$, justifying its formal meaning.
\end{proof}

\subsection{A finite collection of Gamma scales}
For nonzero scales $\rho_1,\ldots,\rho_d$ and constants $c_i$ with
$\sum c_i=1$, put
\[
 K(z)=\sum_{i=1}^dc_iK_{\rho_i}(z),\qquad
 K_\rho(z)=\frac{2\sinh(\rho z/2)}{\rho z}.
\]
Under the digamma substitution, \eqref{ag:eq:profile-density} becomes
\begin{equation}\label{ag:eq:mixed-Gamma}
 \prod_{i=1}^d
 \left(\frac{\Gamma(x+\rho_it/2)}{\Gamma(x-\rho_it/2)}\right)^{c_i/\rho_i}.
\end{equation}
The powers are defined by the analytic logarithm which vanishes at
$t=0$. For positive real scales, real $x>0$, and sufficiently small
real $t$, this is the usual positive real power.

For compactness write
$D_n=\partial_{y^2}^{\,n}$,
$A_{ij}(y)=\cosh((\rho_i-\rho_j)y/2)$, and
$B_{ij}(y)=\cosh((\rho_i+\rho_j)y/2)$.
Theorem~\ref{ag:thm:profile} gives the explicit mixed-scale formula
\begin{align}
 g_{K,r}(y)
 =\frac{(-1)^r}{(r+1)!}\bigg\{&
 \sum_i\frac{c_i^2}{\rho_i^2}D_{r+1}\cosh(\rho_i y)\notag\\
 &+2\sum_{i<j}\frac{c_ic_j}{\rho_i\rho_j}
 \bigl(A_{ij}D_{r+1}B_{ij}-B_{ij}D_{r+1}A_{ij}\bigr)\bigg\}.
 \label{ag:eq:mixed-square}
\end{align}
Indeed, expand the product in \eqref{ag:eq:profile-symbol} and use the
hyperbolic product-to-sum identities before Taylor expanding in $w$.
Coincident scales have a regular limit; no division by
$\rho_i-\rho_j$ appears.

As one concrete check, take $c_1=c_2=1/2$, $\rho_1=1$, and $\rho_2=2$.
The coefficient of $u_2^2$ in the weight-six exponential coefficient,
normalized by $6!$, is
\begin{equation}\label{ag:eq:mixed-example}
 [u_2^2]\,6![t^6]\NF\mathcal C_K(t)=\frac{227}{32}.
\end{equation}
It can also be verified without the symbol: the exponent is
$tu_0+(5/48)t^3u_2+(17/3840)t^5u_4+\cdots$, and
$\NF(u_0u_4)=u_2^2$. The two contributions are $51/16$ and $125/32$.

The general theorem is useful even when $K$ is a polynomial rather
than a Gamma profile. Its proof does not use analytic convergence of
$K$. For convergent profiles the resulting kernels may have additional
analytic representations; the remainder of this article treats the
uniform Gamma profile, whose beta weight gives particularly short
endpoint formulas.
